\documentclass{ltjsbook}
\usepackage{docmute} 
\usepackage{../../myStyle} 

\begin{document}
% \chapter{Rをつかった項目反応理論}
\chapter{Item Response Theory using R}
\label{x11_IRTonR}

% 今回もRを使った演習です。Rを使っての行列計算，Rを使っての因子分析と進めてきましたが，今回はRをつかって項目反応理論を実践的に理解していきましょう。
This time again, it's an exercise using R. We have proceeded with matrix calculations using R and factor analysis using R, but this time let's practically understand the item response theory using R.

% \section{項目反応理論の実際}
\section{Practical Application of Item Response Theory}
% \subsection{データの素描}
\subsection{Overview of the Data}
% 項目反応理論は，テスト理論に関するモデルですから，サンプルデータとして二値データを使いましょう。表\ref{tbl::11_01}に使うデータの一部を示しています\footnote{このデータ全体については，\texttt{IRTsample.csv}というファイル形式で配布していますので，それを読み込んで使ってください}。
The Item Response Theory is a model related to test theory, so let's use binary data as sample data. Some of the data we will be using is shown in Table \ref{tbl::11_01}\footnote{The entire data set is distributed in a file format called \texttt{IRTsample.csv}, so please load and use that}.
\begin{table}[ht]
   \caption{IRTで使うサンプルデータ(一部)}
   \centering
   \begin{tabular}[t]{rrrrrrrrrr}
      \hline
      V1 & V2 & V3 & V4 & V5 & V6 & V7 & V8 & V9 & V10 \\
      \hline
      0  & 0  & 1  & 0  & 0  & 0  & 0  & 0  & 0  & 0   \\
      1  & 1  & 0  & 0  & 0  & 0  & 0  & 0  & 0  & 1   \\
      0  & 1  & 0  & 1  & 0  & 0  & 0  & 0  & 0  & 1   \\
      0  & 1  & 1  & 1  & 1  & 1  & 0  & 0  & 0  & 1   \\
      0  & 1  & 1  & 1  & 1  & 0  & 0  & 0  & 0  & 1   \\
      1  & 1  & 1  & 1  & 0  & 0  & 0  & 0  & 1  & 1   \\
      \hline
   \end{tabular}
   \label{tbl::11_01}
\end{table}


% このように，テスト理論で使うデータは0/1の二値データで，1が正答，0が誤答を表しています。
Thus, the data used in test theory is binary data with 0/1, where 1 represents a correct answer and 0 represents an incorrect answer.

% これらのデータを使ってRで分析をします。IRTを実行するパッケージとして，\texttt{ltm}\citep{ltm}や\texttt{irtoys}\citep{partchev2017package}などがあります。ここでは\texttt{ltm}パッケージの方を使っていきましょう。
We will perform the analysis in R using these data. There are packages to execute IRT, such as \texttt{ltm}\citep{ltm} and \texttt{irtoys}\citep{partchev2017package}. Let's use the \texttt{ltm} package here.

% 分析に先立って，データの記述統計量を確認しておきます。\texttt{code:\ref{code:11_01}}を実行してください。
Before proceeding with the analysis, let's confirm the descriptive statistics of the data. Please execute the code:\ref{code:11_01}.

\begin{lstlisting}[language=c++,label={code:11_01},caption={データの読み込みから記述統計量の計算まで}]
    rm(list=ls())
    library(tidyverse)
    library(ltm)
    dat <- read_csv("IRTsample.csv")
    ltm::descript(dat)
\end{lstlisting}


% \paragraph{コード解説}
\paragraph{Code Explanation}
\begin{description}
   \item[1行目] 環境の初期化
   \item[2-3行目] パッケージの読み込み
   \item[4-5行目] サンプルデータ\texttt{IRTsample.csv}を読み込む。文字化けする人は，\verb|read_csv|関数のオプションとして\texttt{locale=locale(encoding="UTF-8")}を追加してください。
   \item[6行目] ltmパッケージに含まれる\texttt{descript}関数を実行
\end{description}


% いろいろな情報が一気に出力されるので，順にみていきます。まず\texttt{Rの出力\ref{RS:out11_01}}のところです。ここには各変数の0と1の比率が入っています。この例では，\texttt{V1}は0が71.0\%，1が29.0\%含まれているということです。最後の\texttt{logit}は，0から1の間に入る数字$p$にたいして，$logit(p)=\log\left(\frac{p}{1-p}\right)$で計算される量です。ここでは$p$は正答率であり，この\texttt{logit}の値が負であれば0の比率が，正であれば1の比率が多かった，ということを意味します。0/1がちょうど半々であれば，\texttt{logit}の値は0になる，そういう数字です。これを見るだけでも，\texttt{V8}はずいぶんと正答率が低かったんだな，\texttt{V3}は逆に正答率が高かったんだな，ということがわかります。
Various information is output all at once, so let's go through them in order. First of all, there is the section \texttt{R's output\ref{RS:out11_01}}. Here, you can find the ratio of 0 and 1 for each variable. In this example, \texttt{V1} contains 71.0% for 0 and 29.0% for 1. The last \texttt{logit} amount is calculated by $logit(p) = \log\left(\frac{p}{1-p}\right)$ for a number $p$ that falls between 0 and 1. Here, $p$ is the correct answer rate, and if this \texttt{logit} value is negative, it means that the ratio of 0 was higher, and if positive, it means that the ratio of 1 was higher. If 0/1 is exactly half, the \texttt{logit} value will be zero, such a number. Just by looking at this, you can see that \texttt{V8} had a particularly low correct answer rate, and \texttt{V3} had a high correct answer rate.
\begin{Rscreen}{記述統計の出力1}{out11_01}
	\begin{verbatim}
    Proportions for each level of response:
    0     1   logit
V1  0.710 0.290 -0.8954
V2  0.114 0.886  2.0505
V3  0.186 0.814  1.4762
V4  0.206 0.794  1.3492
V5  0.670 0.330 -0.7082
V6  0.832 0.168 -1.5999
V7  0.666 0.334 -0.6901
V8  0.884 0.116 -2.0309
V9  0.652 0.348 -0.6278
V10 0.254 0.746  1.0774
\end{verbatim}
\end{Rscreen}


% 次のセクション\texttt{Rの出力\ref{RS:out11_02}}には，正答数の分布が示されています。10問全部正解したのは3人だけ，6問正解した人が104人だった，といったことがわかります。
The next section, \texttt{R's Output\ref{RS:out11_02}}, shows the distribution of correct answers. We can see that only 3 people answered all 10 questions correctly, and 104 people answered 6 questions correctly.
\begin{Rscreen}{記述統計の出力2}{out11_02}
   \begin{verbatim}
        Frequencies of total scores:
        0  1  2  3  4   5  6  7  8  9 10
   Freq 4 19 41 62 97 104 66 54 36 14  3
    \end{verbatim}
\end{Rscreen}


% 次のセクション\texttt{Rの出力\ref{RS:out11_03}}には，\keyterm{点双列相関係数}{てんそうれつそうかんけいすう}{Point Biserial correlation}が計算されています。この相関係数は2値と連続値の相関係数であり，ここでは各変数(二値)と合計得点(0から10)の相関係数を計算しています。この合計得点の計算方法には，その変数自身を含める場合(Included)と含めない場合(Excluded)の二種類があります。具体的に説明したほうがわかりやすいでしょう。たとえば変数\texttt{V1}についていうと，「その変数自身を含める場合」の合計得点とは$V1+V2+V3+V4+V5+V6+V7+V8+V9+V10$で計算しています。「含めない場合」は$V2+V3+\cdots+V10$とするわけです。含める場合の相関係数が0.5207，含めない場合の相関係数が0.2480です。当然自身が含まれているほうが合計得点に深く関わっているので相関係数が高くなります。裏を返せば，含める場合と含めない場合の相関係数に大きな差があるとき，その変数はテスト全体に大きく貢献する重要な変数だといえるでしょう。
The next section, titled \texttt{R's output\ref{RS:out11_03}}, calculates the \keyterm{Point Biserial correlation coefficient}{}. This correlation coefficient is a correlation coefficient between binary and continuous values. Here, we calculate correlation coefficients for each variable (binary) and total scores (from 0 to 10). There are two types of total score calculation methods: 'Included', which includes the variable itself, and 'Excluded', which does not.

Let me explain more specifically. For example, for the variable \texttt{V1}, the total score 'including the variable itself' is calculated as $V1+V2+V3+V4+V5+V6+V7+V8+V9+V10$, whereas 'excluding the variable itself' is $V2+V3+\cdots+V10$. The correlation coefficient for the 'included' case is 0.5207, while that for the 'excluded' case is 0.2480. Naturally, the correlation coefficient is higher when the variable itself is included, as it is deeply involved in the total score. Conversely, when there is a significant difference between the correlation coefficients for the 'included' and 'excluded' cases, it can be said that the variable makes a significant contribution to the overall test.
\begin{Rscreen}{記述統計の出力3}{out11_03}
   \begin{verbatim}
Point Biserial correlation with Total Score:
    Included Excluded
V1    0.4228   0.2118
V2    0.3959   0.2502
V3    0.5202   0.3567
V4    0.5209   0.3501
V5    0.5449   0.3470
V6    0.4635   0.2983
V7    0.5440   0.3452
V8    0.2901   0.1350
V9    0.4992   0.2890
V10   0.5890   0.4180
    \end{verbatim}
\end{Rscreen}


% 次のセクション\texttt{Rの出力\ref{RS:out11_04}}には，\keyterm{$\alpha$係数}{あるふぁけいすう}{alpha coefficient}が算出されています。これも全体の場合と，特定の変数を抜いた場合とを計算しています。
In the next section, R Output\ref{RS:out11_04}, the alpha coefficient is calculated. This is done both for the overall case and for the case where specific variables are excluded.
% 心理尺度における信頼性係数としての$\alpha$係数は，0.8以上であれば信頼性がある，といった目安がありますが，二値データの場合はここに示したようにそれよりもずいぶんと低くなります。これは二値データの分散が小さくなることから仕方のないことですし，項目反応理論で分析する上での信頼性は情報量関数として表されますから，問題ではありません。
In psychological scales, there is a guideline that if the $\alpha$ coefficient of reliability is 0.8 or higher, it is considered reliable, but in the case of binary data, it becomes considerably lower as shown here. This is inevitable because the variance of binary data is small, and since the reliability in the analysis using item response theory is represented as an information function, it is not a problem.
\begin{Rscreen}{記述統計の出力4}{out11_04}
   \begin{verbatim}
        Cronbach's alpha:
        value
All Items     0.6341
Excluding V1  0.6303
Excluding V2  0.6195
Excluding V3  0.5978
Excluding V4  0.5986
Excluding V5  0.5982
Excluding V6  0.6100
Excluding V7  0.5987
Excluding V8  0.6380
Excluding V9  0.6130
Excluding V10 0.5817
    \end{verbatim}
\end{Rscreen}


% 最後にセクション\texttt{Rの出力\ref{RS:out11_05}}では，変数同士の関連がとくに弱いものが示されています。変数同士のペアについて度数分布表を作成し，$\chi^2$検定によって\texttt{p.value}を算出しています。この値が小さければ，2つの変数の関係が強いことを示していると解釈できるのですが，ここにあげているように大きな数字になっているということは，変数間の関係が弱いと読み取ることができます。
In the final section, \texttt{R Output\ref{RS:out11_05}}, it shows that the relationships between variables are particularly weak. Frequency distribution tables are created for each pair of variables, and the \texttt{p.value} is calculated using a $\chi^2$ test. If this value is small, it can be interpreted that the relationship between the two variables is strong, but as shown here, if it turns out to be a large number, it can be read as the relationship between variables is weak.
\begin{Rscreen}{記述統計の出力4}{out11_05}
	\begin{verbatim}
Pairwise Associations:
   Item i Item j p.value
1       2      8   0.625
2       1      8   0.409
3       1      2   0.347
4       6      8   0.303
5       4      8   0.234
6       7      8   0.222
7       5      8   0.195
8       1      4   0.191
9       1      3   0.166
10      3      8   0.124
    \end{verbatim}
\end{Rscreen}

% \subsection{IRTモデルの実際}
\subsection{Practical Application of the IRT Model}
% データの感じがわかったところで，IRTモデルを実行していきましょう。
Now that we understand the feel of the data, let's proceed with executing the IRT model.

% 今回は3つのモデルを試してみます。確認のために，少しモデル式に言及しておきます。困難度母数だけに注目するのが\keyterm{1パラメータ・ロジスティックモデル}{1ぱらめーたろじすてぃっくもでる}{1 parameter logistic model, 1PL}で，モデル式は次の通りです。
This time, we will try three models. For confirmation, let's mention a bit about the model equation. The \keyterm{1 parameter logistic model, 1PL} that focuses only on the difficulty parameter, and the model equation is as follows.

\[ p_j(\theta) = \frac{1}{1+exp(-1.7 (\theta -b_j))} \]

ここで$p_j(\theta)$は能力が$\theta$のひとが項目$j$に正答する確率であり，$b_j$が\keytermL{困難度母数}{こんなんどぼすう}{difficulty parameter}と呼ばれます。これは開発者の名前を使って，別名\keyterm{ラッシュモデル}{らっしゅもでる}{rasch model}とも呼ばれています。

項目の識別力にも注目するのが\keyterm{2パラメータ・ロジスティックモデル}{2ぱらめーたろじすてぃっくもでる}{2 parameter logistic model, 2PL}で，モデル式は次の通りです。

\[ p_j(\theta) = \frac{1}{1+exp(-1.7 a_j(\theta -b_j))} \]

% ここで$a_j$は\keyterm{識別力母数}{しきべつりょくぼすう}{descriminant paramter}と呼ばれます。
Here, $a_j$ is called the discriminant parameter.

% 最後に\keyterm{3パラメータ・ロジスティックモデル}{3ぱらめーたろじすてぃっくもでる}{3 parameter logistic model, 3PL}では，困難度，識別力に加えて\keyterm{当て推量母数}{あてすいりょうぼすう}{guessing parameter}というのを推定します。モデル式は次の通りです。
Finally, in the \keyterm{3 parameter logistic model, 3PL}{3ぱらめーたろじすてぃっくもでる}{3 parameter logistic model, 3PL}, in addition to difficulty and discriminatory power, a \keyterm{guessing parameter}{あてすいりょうぼすう}{guessing parameter} is estimated. The model formula is as follows.

\[ p_j(\theta) =c+ \frac{1-c}{1+exp(-1.7 a_j(\theta -b_j))} \]

もちろん式ではピンと来ないという人もいると思いますので，実際にどのように表されるかみていきましょう。
\subsubsection{1PLモデルの場合}

\texttt{ltm}パッケージを使って1PLモデルを実行します。\texttt{code:\ref{code:11_02}}を実行してください。

\begin{lstlisting}[language=c++,label={code:11_02},caption={1PLモデルの実行}]
result.1pl <- rasch(dat)
print(result.1pl)
plot(result.1pl,type="ICC")
plot(result.1pl,type="IIC")
plot(result.1pl,type="IIC",items = 0)
\end{lstlisting}

\paragraph{コード解説}
\begin{description}
	\item[1行目] 1PLモデル(ラッシュモデル)の実行
	\item[2行目] 結果の出力
	\item[3行目] \keyterm{項目特性曲線}{こうもくとくせいきょくせん}{Item Characteristic Curve}のプロット。\texttt{plot}関数のオプションで表示する情報をICCと指定。
	\item[4行目] \keyterm{項目情報曲線}{こうもくじょうほうきょくせん}{Item Information Curve}のプロット。\texttt{plot}関数のオプションで表示する情報をIICと指定。
	\item[5行目] \keyterm{テスト情報曲線}{てすとじょうほうきょくせん}{Test Information Curve}のプロット。\texttt{plot}関数のオプションで表示する項目番号をゼロにするとTICになる。
\end{description}

2行目で結果の出力を行っていますが，出力\ref{RS:out11_06}のように表されます。
\begin{Rscreen}{2PLモデルの結果出力}{out11_06}
	\begin{verbatim}
Call:
rasch(data = dat)

Coefficients:
 Dffclt.V1   Dffclt.V2   Dffclt.V3   Dffclt.V4   Dffclt.V5   Dffclt.V6   
     0.992      -2.208      -1.615      -1.480       0.787       1.744
 Dffclt.V7   Dffclt.V8   Dffclt.V9  Dffclt.V10      Dscrmn  
     0.768       2.187       0.699      -1.189       1.126  
Log.Lik: -2468.743
\end{verbatim}
\end{Rscreen}
1PLモデルで表現する項目の違いは\keytermL{困難度母数}{こんなんどぼすう}だけですから，この数字(\texttt{Dffclt})が小されければ簡単な項目，大きければ難しい項目だったことがわかります。たとえば\texttt{V2}の困難度母数は$b_2=-2.208$ですが，記述統計のところで見た正答率(88.6\%)を考えると，簡単な問題だったことがわかりますね。\keyterm{識別力母数}{しきべつりょくぼすう}は\texttt{Dscrmn}として1.126と推定されています。これは項目を通じて一定です。このことは図\ref{fig::11_01}の左にあるICCをみて，すべての曲線の傾きが同じことで確認できます。
\begin{figure}[ht]
	\centering
	\includegraphics[keepaspectratio,width=0.8\linewidth]{../images/11_IRTonR/Rplot11_01.png}
	\caption{1PLのさまざまなプロット。左からICC,IIC,TIC}
	\label{fig::11_01}
\end{figure}

図\ref{fig::11_01}の左が数値的出力結果を図示したものです。これを見た方がわかりやすいかもしれません。x軸は$\theta$，つまり測定する潜在変数の値になっており，テストで言えばある人の学力$\theta_i$があったときに，各問いにどれぐらいの確率で正答するかを表していることになります。これをみると\texttt{V2}がもっとも簡単で，\texttt{V8}が最も難しかったテストであることがわかります。

また図\ref{fig::11_01}の中央にあるのがIICで，それぞれの項目がどのあたりの能力値$\theta$を測定するときにもっとも敏感になるかを表しています。これは項目反応理論でいう\keytermL{信頼性}{しんらいせい}の表現であることを思い出してください。最後に図の左側にあるのは，このIICを10問分足し合わせたテスト全体の情報曲線，TICになります。これをみると，$\theta=0$よりやや右側のあたりにピークがきているので，このテスト全体としては平均より高い学力の人を測定するときに，最も鋭敏に働くことがわかるでしょう。

\subsubsection{2PLモデルの場合}

ひきつづき\texttt{ltm}パッケージを使って，2PLモデルを実行していきましょう。\texttt{code:\ref{code:11_03}}を実行してください。

\begin{lstlisting}[language=c++,label={code:11_03},caption={2PLモデルの実行}]
result.2pl <- ltm(dat~z1)
print(result.2pl)
plot(result.2pl,type="ICC")
plot(result.2pl,type="IIC")
plot(result.2pl,type="IIC",items = 0)
\end{lstlisting}

\paragraph{コード解説}
\begin{description}
	\item[1行目] 1PLモデル(ラッシュモデル)の実行
	\item[2行目] 結果の出力
	\item[3行目] \keyterm{項目特性曲線}{こうもくとくせいきょくせん}{Item Characteristic Curve}のプロット
	\item[4行目] \keyterm{項目情報曲線}{こうもくじょうほうきょくせん}{Item Information Curve}のプロット
	\item[5行目] \keyterm{テスト情報曲線}{てすとじょうほうきょくせん}{Test Information Curve}のプロット
\end{description}

2行目で結果の出力を行っていますが，出力\ref{RS:out11_15}のように表されます。
\begin{Rscreen}{2PLモデルの結果出力}{out11_15}
	\begin{verbatim}
Call:
ltm(formula = dat ~ z1)

Coefficients:
     Dffclt  Dscrmn
V1    1.602   0.603
V2   -2.078   1.228
V3   -1.213   1.892
V4   -1.199   1.620
V5    0.759   1.176
V6    1.698   1.175
V7    0.740   1.174
V8    4.133   0.516
V9    0.866   0.828
V10  -0.880   2.020

Log.Lik: -2446.104
    \end{verbatim}
\end{Rscreen}

2PLモデルで表現する項目の違いは，\keytermL{困難度母数}{こんなんどぼすう}と\keytermL{識別力母数}{しきべつりょくぼすう}の2つです。識別力母数を入れることで，IIC曲線の傾きもデータに合わせて調整できるようになりました。識別力母数のデフォルトは(推定しないのであれば)1.0ですから，ここから大きく逸脱するような項目には注意です。今回は\texttt{V8}をみると，困難度が非常に高いこともわかりますが，識別力が非常に低くなっているので，この項目に誤答したからと言ってすぐさま成績が悪い，と判断できるほどではないことがわかります。

最後の一行に\texttt{Log.Lik}とありますが，これは\keyterm{対数尤度}{たいすうゆうど}{log likelihood}を表しています。これが大きいほど，データとモデルの適合度が高いことを意味します。1PLのときの\texttt{Log.Lik}が $-2468.743$でしたから，1PLモデルより2PLモデルの方が当てはまりがよかったと言えるでしょう。

プロットによる出力も見ておきましょう(図\ref{fig::11_02})。識別力母数が入ったことで，ICCの傾きが項目によってさまざまになりますし，IICやTICの見た目もすっかり変わってしまいましたが，読み取り方は同じです。
\begin{figure}[ht]
	\centering
	\includegraphics[keepaspectratio,width=0.8\linewidth]{../images/11_IRTonR/Rplot11_02.png}
	\caption{2PLのさまざまなプロット。左からICC,IIC,TIC}
	\label{fig::11_02}
\end{figure}


\subsubsection{3PLモデルの場合}

最後に\texttt{ltm}パッケージを使って，3PLモデルを実行していきましょう。\texttt{code:\ref{code:11_04}}を実行してください。

\begin{lstlisting}[language=c++,label={code:11_04},caption={3PLモデルの実行}]
result.3pl <- tpm(dat)
print(result.3pl)
plot(result.3pl,type="ICC")
plot(result.3pl,type="IIC")
plot(result.3pl,type="IIC",items = 0)
\end{lstlisting}


コードは一行目の関数以外ほとんど同じですので，逐一解説は致しませんが，出力のところを見ておきましょう(出力\ref{RS:out11_07})。
\begin{Rscreen}{3PLモデルの結果出力}{out11_07}
	\begin{verbatim}
Call:
tpm(data = dat)

Coefficients:
     Gussng  Dffclt  Dscrmn
V1    0.230   1.407  15.123
V2    0.638  -0.610   2.232
V3    0.225  -0.825   2.907
V4    0.324  -0.637   2.324
V5    0.149   0.931   2.992
V6    0.043   1.613   1.706
V7    0.074   0.879   1.509
V8    0.000   4.490   0.471
V9    0.133   1.144   1.255
V10   0.278  -0.438   3.085

Log.Lik: -2431.371
        
    \end{verbatim}
\end{Rscreen}

3PLモデルは困難度，識別力に加えて\keytermL{当て推量母数}{あてすいりょうぼすう}が計算されます。これはICCの形(図\ref{fig::11_03}の左)を見ればわかりますが，曲線全体が底上げされたようになっているものがあります。この最低ラインの高さが当て推量母数の大きさです。図や数値から，\texttt{V2}のそれが大きい数字であることが読み取れます。ICCは$\theta$の関数である正答率ですが，これが$\theta$の値にかかわらず一定の値を取るということは，どれだけ能力が低くてもその確率で正答してしまうことを意味しています。なので「適当に推測して答えを当てられる大きさ」という意味で当て推量母数と呼ばれているのです。

\begin{figure}[ht]
	\centering
	\includegraphics[keepaspectratio,width=0.8\linewidth]{../images/11_IRTonR/Rplot11_03.png}
	\caption{3PLのさまざまなプロット。左からICC,IIC,TIC}
	\label{fig::11_03}
\end{figure}

今回当て推量母数を含めて考えたことで，IICやTICがまた大きく変化しました。この数字が高すぎるV2は，テストとしてはあまり良い項目ではなかったのかもしれません。また，モデルの適合度としては3PLモデルが最もよかったのですが，2PLモデルとそれほど大きく違うわけでもありません。
どのモデルでどのように推定するべきか，推定されたモデルはどういう仮定を持っているのかについて，しっかり理解した上で利用するようにしましょう。

\section{段階反応モデルの実際}

続いて\keytermL{段階反応モデル}{だんかいはんのうもでる}の練習に進みましょう。項目反応理論は正答/誤答の二値データに対して用いるものでしたが，これを多段階反応に拡張したモデルの1つが段階反応モデルです。多段階反応とは，「まったく当てはまらない」を1，「非常に当てはまる」を7とする，といったような\keytermL{リッカート尺度}{りっかーとしゃくど}のようなスタイルの調査項目です。反応させるカテゴリの段階数に応じて，5件法とか7件法と呼ばれたりします。

リッカート尺度の場合は，シグマ法をつかって累積度数からカテゴリに該当する値を算出し，それを元に態度得点を作るという方法でした。とはいえこのやり方でやっても，いい尺度であればそのまま素点を尺度値としてもほとんど問題がないため，シグマ法が回りくどい方法だと思われたのか，使われなくなってしまいました。しかし実際には「いい尺度であれば」という条件を満たしているかどうか，しっかり確認しなければなりません。歪んだ分布や偏ったカテゴリ反応なのに，素点をそのまま尺度値とするのは不適切なことがあります。また，リッカート尺度で撮られたデータは因子分析をすることによって，解釈度に分類されて，因子ごとのスコアを算出して利用されることが一般的です。しかし，尺度値が適切でなければ，ピアソンの積率相関係数も適切ではなく，そこから計算が始まる因子分析のモデルにも疑義が生じます。

カテゴリ反応はあくまでも並びの問題であって，\keytermL{順序尺度水準}{じゅんじょしゃくどすいじゅん}の情報しか持たないと考えるのであれば，その背後にある態度のような連続体を仮定し，この態度$\theta$が大きくなることで徐々に反応カテゴリが変わる，という\keytermL{段階反応モデル}{だんかいはんのうもでる}のほうが適切な分析方法になるわけです。

項目反応理論(の段階反応モデル)を使って確認することは決して難しくありません。今回は因子分析の時に使った\texttt{bfi}データの一部を使って段階反応モデルを実践してみましょう。一部を使うとしたのは，段階反応モデルは項目反応理論の中間ですから，潜在変数が1次元(一因子)である，という仮定があるからです。ここでは\texttt{N}因子(Neuroticism，神経質さ)を使った例を示しています。

\begin{lstlisting}[language=c++,label={code:11_05},caption={GRMの実行}]
library(psych)
dat <- bfi %>%
  dplyr::select(-gender, -education, -age) %>%
  dplyr::select(starts_with("N"))
result.grm <- grm(dat)
result.grm
plot(result.grm, type = "ICC", items = 1)
plot(result.grm, type = "IIC", items = 1)
plot(result.grm, type = "ICC", items = 4)
plot(result.grm, type = "IIC", items = 4)
plot(result.grm, type = "IIC", items = 0)
\end{lstlisting}

\paragraph{コード解説}
\begin{description}
	\item[1行目] \texttt{psych}パッケージを読み込む。サンプルデータ\texttt{bfi}を用いるためです。
	\item[2-4行目] データの加工。\texttt{bfi}データから，性別，最終学歴，年齢などの情報を取り除き，またNで始まる変数だけ選出(\texttt{select})しています。
	\item[5行目] 段階反応モデルの実行。
	\item[6行目] 結果の出力。
	\item[7行目] 項目1についてのICC，厳密には\keyterm{項目反応カテゴリ特性曲線}{こうもくはんのうかてごりとくせいきょくせん}{Item Response Category Characteristic Curve}と呼びます。
	\item[8行目] 項目1についてのIICを表示させています。
	\item[9行目] 項目4についてのICCを表示させています。
	\item[10行目] 項目4についてのICCを表示させています。
	\item[11行目]項目群についてのTICを表示させています。
\end{description}

\texttt{code:\ref{code:11_05}}の6行目で，分析結果の数値的出力をさせていますので，それをみておきましょう(出力\ref{RS:out11_25})。

ここで，\texttt{Dscrmn}とあるのは識別力母数ですが，困難度母数に対応するのが\texttt{Extrmt1}から\texttt{Extrmt5}にあたります。この尺度は6件法ですので，反応カテゴリが変わる切れ目，すなわち\keytermL{閾値}{いきち}が5つあるのです。これらが2PLモデルでいうところの困難度に対応すると言ってもいいでしょう。

\begin{Rscreen}{GRMモデルの結果出力}{out11_25}
	\begin{verbatim}
Call:
grm(data = dat)

Coefficients:
    Extrmt1  Extrmt2  Extrmt3  Extrmt4  Extrmt5  Dscrmn
N1   -0.810   -0.093    0.342    0.985    1.720   3.125
N2   -1.366   -0.555   -0.111    0.648    1.483   2.890
N3   -1.187   -0.298    0.123    0.876    1.767   2.026
N4   -1.564   -0.355    0.238    1.240    2.280   1.277
N5   -1.296   -0.125    0.494    1.479    2.530   1.112

Log.Lik: -21721.91
\end{verbatim}
\end{Rscreen}

数値を見ていてもピンとこないかもしれませんので，図でこの各項目の特徴を確認しておきましょう。変数\texttt{N1}と\texttt{N4}のIRCCCを図\ref{fig::11_04}に示しました。この図に表されている複数の曲線が，$\theta$に対応した各カテゴリに反応する確率を表しています。図の左，\texttt{N1}は綺麗なカーブが描かれており，各カテゴリのピークが見て取れますが，\texttt{N4}は反応カテゴリ3の山が埋もれてしまっていることがみてとれます。この項目については，6段階ではなく5段階で反応を求めるのがよかったのかもしれません\footnote{ちなみに\texttt{N1}は「すぐ怒る(Get angry easily)」，\texttt{N4}は「割と凹む(Often feel blue)」です。日本語訳は小杉の意訳で定訳ではありません。あしからず。}。

\begin{figure}[ht]
	\centering
	\includegraphics[keepaspectratio,width=0.8\linewidth]{../images/11_IRTonR/Rplot11_04.png}
	\caption{GRMのIRCCC}
	\label{fig::11_04}
\end{figure}

原理的には難しいことのように感じますが，コードを書いて実践する分にはほとんど苦労がなく，二値モデルと同じように考えることができます。コードにはIICやTICを描くものもありますが，1PL,2PL,3PLモデルと変わらない感じで多段階モデルが実行できますね。
みなさんも心理尺度を使う時は，盲目的にスコアリングし，因子分析するのではなく，どういう仮定の元でどういう計算をしているかを理解しながら，データや調査協力者にとって最適な分析をするように心がけてください。

\section{カテゴリカル因子分析との対応}
さて多段階反応による分析も簡単にできることがわかりましたが，気になるのは段階反応モデルが一因子構造しか仮定していないところです。bfiデータはビッグファイブと呼ばれるように五因子(5次元)の仮定があるのですが，一因子ずつ\texttt{grm}を実行し，あとで統合するのは大変です。

でも大丈夫。簡単に多次元に展開する方法があります。
いわゆる一般的な因子分析は，ピアソンの積率相関係数を元に(積率相関係数の相関行列$\bm{R}$の固有値分解から)因子を求めていくのですが，この積率相関係数が\keytermL{間隔尺度水準}{かんかくしゃくどすいじゅん}を仮定したもので，シグマ法で計算された尺度値を使っていないのであれば疑義が残る，という話でした。
この元になる相関係数が，潜在的連続体を仮定した順序尺度水準向けの相関係数である，\keytermL{ポリコリック相関係数}{ぽりこりっくそうかんけいすう}であれば，その問題は解決します。反応段階が順序尺度を仮定した因子分析のことをとくに\keytermL{カテゴリカル因子分析}{かてごりかるいんしぶんせき}と言います。カテゴリカル因子部な席は，前回紹介した\texttt{psych}パッケージに少しオプションを足すだけで簡単に実行できるものです。

\begin{lstlisting}[language=c++,label={code:11_06},caption={GRMの実行}]
dat <- bfi %>% dplyr::select(-gender, -education, -age)
result.fa <- fa.poly(dat, nfactors = 5, rotate = "geominQ")
print(result.fa, sort = T, cut = 0.3)
\end{lstlisting}

\paragraph{コード解説}
\begin{description}
	\item[1行目] サンプルデータ\texttt{bfi}の加工をしています。
	\item[2行目] ポリ個リック相関係数に基づく因子分析の実行
	\item[3行目] 結果の出力。
\end{description}

コードの2行目を見ていただくとわかる通り，ただの\texttt{fa}関数ではなく\texttt{fa.poly}としただけで，あとは因子分析と同じです。この\texttt{.poly}のところがポリコリック相関係数を使うよう指定しているところで，これで実質的に測定尺度水準が順序尺度水準であると想定した因子分析をしたことになっています。

カテゴリカル因子分析は，段階反応モデルと数学的に等価であることが知られています。出自が違いますので，分析の数字や出力方法に違いがありますが\footnote{たとえば因子分析の文脈では，因子負荷量を提示して単純構造を目指すのに対し，項目反応理論の文脈ではIRCCCを描いたりTICを描いたりして各反応カテゴリの特徴を精査します。}，実質的には同じことをしていると理解してください。また，Rではたった数文字追加するだけで適切な分析方法になるのですから，使わない手はないと思います。

カテゴリカル因子分析は，因子分析の文脈から順序尺度水準の項目を扱えるように，と発展してきたものです，これに対して，項目反応理論の文脈から心理尺度を扱えるように，と発展してきたルートもあります。これは\keyterm{多次元項目反応理論}{たじげんこうもくはんのうりろん}{multi-dimensional item response theory}と呼ばれるもので，これを提供するRパッケージもあります。それが\texttt{mirt}パッケージ\citep{mirt}です。最後にこれを使ったコードの例を示しておきます。この計算には時間がかかりますが，完全情報最尤推定によって項目反応理論で算出するような，精度の高い\keytermL{因子得点}{いんしとくてん}を得ることができるのは利点だと言えるでしょう。

\begin{lstlisting}[language=c++,label={code:11_06b},caption={mirtの実行}]
library(mirt)
result.mirt <- mirt(dat, 5)
print(result.mirt)
result.mirt %>% summary(rotate = "geominQ")
\end{lstlisting}


\section{課題}
今日の授業でおこなったすべての次の計算をするRコードを提出してください。ファイル形式はRスクリプトかRmdとします。なお提出されたコード単体でバグがなく動くことが確認できないものは，未提出扱いになります。コードの書き方などわからないところがあれば，曜日別TAか小杉までメールで連絡し，指導を受けてください。


\end{document}
